% Conclusiones del informe final. Las cifras provienen de la sección de resultados.
\section{Conclusiones}
\label{sec:conclusiones}

El proyecto partió de una pregunta: si un modelo aditivo legible alcanza en COFINFAD la precisión de los
ensambles de árboles, y a qué costo. La respuesta de la sección~\ref{sec:resultados} es que la alcanza, no la
supera. La EBM obtiene el mayor AUC de prueba de los ocho modelos de clasificación (0,8289) y el segundo RMSE de
regresión (0,05925, a 0,00008 de XGBoost); supera con significación estadística, en todas las pruebas y tras la
corrección de Holm, a los modelos lineales, a k-NN y a Naive Bayes (cuatro de las siete referencias de cada
tarea); y frente a XGBoost, Random Forest y el árbol de decisión ninguna prueba al 5~\% encuentra diferencia:
la probabilidad de equivalencia práctica va de 0,94 a 0,997 en clasificación y es 0,9998 con XGBoost en
regresión, y solo el \emph{Model Confidence Set} al 90~\% detecta una ventaja sistemática de XGBoost de
décimas de milésima de RMSE. Frente a la entrega anterior, en su misma partición y etiqueta, la EBM sube el AUC
de 0,6832 a 0,7187 y baja el RMSE de 0,06166 a 0,05926 ($p < 0{,}001$), y cubre el 98~\% del techo de
información que deja la exclusión de \ap: el empate entre los modelos flexibles es una propiedad de los datos,
no una debilidad del modelo.

El estudio de ablación cerró sus cinco hipótesis. Se confirmó que la aditividad basta (sin interacciones se
pierden 0,0002 de AUC, $p = 0{,}91$), que el balanceo no cambia el AUC y solo desplaza el punto de operación (el
\emph{recall} con umbral 0,5 pasa de 0,355 a 0,707 con ponderación, al mismo AUC), que la búsqueda de baja
fidelidad conserva el orden de las configuraciones y ahorra el 66~\% del tiempo, y que la eficiencia por
evaluación sí distingue a los optimizadores (la bayesiana tiene el menor arrepentimiento) aunque la métrica
final solo separe a Random Search en clasificación. Se refutó, en cambio, la hipótesis de que la variabilidad
por semilla fuera del orden de las diferencias entre modelos: la EBM es casi determinista (desviación de
0,00004 de AUC), más estable que Random Forest y XGBoost. La sensibilidad a los hiperparámetros es baja (AUC
entre 0,8267 y 0,8292 al variarlos uno a uno) y la validación anidada predijo la prueba casi sin sesgo
(correlación 0,997 en 124 corridas), de modo que las conclusiones no dependen de una configuración afortunada.

Dos lecciones son prácticas. La primera es que la significancia estadística no es relevancia práctica: con la
capacidad de campaña del 25~\% todos los modelos salvo Naive Bayes identifican casi a los mismos clientes (la
EBM 1.228 de los 2.399 de riesgo alto y la regresión logística 1.225), aunque el AUC los separe con
$p < 0{,}001$; la ventaja de los modelos flexibles aparece en otras capacidades (con el 50~\% capturan el 91~\%
de la fuga frente al 77~\% de la logística), así que el modelo debe elegirse con la capacidad de la campaña en
mente y con las probabilidades recalibradas (el ECE baja de 0,170 a 0,003--0,010 con regresión isotónica). La
segunda es el costo: la EBM es el modelo más lento de ajustar (26,7~s, de 15 a 17 veces XGBoost), pero su
inferencia sobre los 9.745 clientes de prueba tarda 0,067~s, como XGBoost, y ocupa 1,5~MB; para una campaña
que ajusta el modelo una vez y lo consulta sobre toda la cartera, el costo relevante es el de inferencia. Lo que
el ajuste más caro compra es la legibilidad: cuatro variables reúnen el 82~\% de la importancia y sus funciones
de forma muestran directamente la no linealidad que los modelos lineales no captan (la edad no monótona, la
meseta de la frecuencia de uso, los escalones de la satisfacción), sin explicador auxiliar. El análisis de
errores localizó los fallos: la cohorte de 59--60 años, con el 46~\% de riesgo alto, es la más difícil, y los
errores grandes de la regresión siguen una escalera por número de productos activos, es decir, son la parte del
objetivo que depende de la variable excluida y que ningún modelo puede recuperar sin ella.

Las limitaciones acotan estas conclusiones. Hay un único conjunto de datos, con rasgos sintéticos (tres
cohortes de edad, una escalera de productos activos, transacciones sin señal), y el objetivo es el puntaje de
otro modelo, no la fuga observada, de modo que lo aprendido es a reproducir ese puntaje. Con cinco pliegues la
$t$ corregida tiene cuatro grados de libertad y poca potencia, McNemar en el punto de operación es insensible a
las diferencias que el AUC sí ve, el SVM es lineal, la regresión no incluye un modelo bayesiano, los tiempos son
de un solo equipo, y la $t$ corregida, McNemar, el análisis bayesiano y los estudios de sensibilidad se
añadieron después de ver la prueba y se declaran exploratorios. Como trabajo futuro se propone validar en el
tiempo con datos reales de fuga observada (entrenar con un periodo y probar con el siguiente), incorporar el
costo de cada acción de retención para fijar el umbral y la capacidad, acelerar el ajuste de la EBM con una
parada temprana más exigente o restringiendo los predictores a los que tienen señal, aumentar la potencia de las
pruebas pareadas con validación cruzada repetida, y explotar las secuencias de transacciones con modelos que las
traten como series.
